% ===================================================================
%  8. TAULA — Uzta. — Ehiza, Mahats-bilketa, Arrantza.
%  Traduction 2026 d'apres texto/io, controlee sur texto/fr, texto/en
%  et texto/es. Consigne : voir texto/eu/10-taula-01.tex.
%
%  L'appel de note est dans le TITRE de la scene — « Uzta (*) ».
%
%  LE BLOC c2-03-1 EST UN ECART DECLARE dans renvoji.py, comme dans
%  toutes les colonnes.
% ===================================================================

%%K t08-noto-1 noto
\VUnotes{143.2mm}{%
\textit{(*) Hitz hau zehatza ez delako: lihoaren uzta, mahats-bilketa, sagar-bilketa?
Agian hobe zatekeen «} \VUgras{mesono} \textit{» hartzea, } Mondo \textit{XIV,
242. or.-n proposatua. Baina oraingoz erro berri hori ez du Akademiak onartu eta
ido-mugimenduaren ohiko arauen araberako eztabaidaren zain dago.}%
}

%%K t08-apar-1 apar
\VUsaut{5.0mm}
\VUtitre{15.0pt}{60}{8. TAULA}
\VUsaut{3.0mm}
\VUcentre{13.2pt}{90}{\textit{Lehen agerraldia.}}
\VUsaut{3.0mm}
\VUpk{10.2pt}{40}{Uzta (*).}
\VUsaut{4.0mm}
\VUpk{10.2pt}{40}{Landen aurpegia.}
\VUsaut{3.0mm}

%%K t08-c1-01-1 p
1. --- \VUgras{Uda} etorrita, landek berriro aldatu dute aurpegia. Ildoari
emandako hazia erne da; lurretik landare berde txiki bat altxatu da eta bere
galburua eman du. Alea bete da, gero heldu, eta orain uzta biltzea baino ez da
geratzen.

%%K t08-c1-02-1 p
2. --- \VUgras{Landetako loreak} zabaldu dira. \VUgras{Anabelarrak}
\textsuperscript{(1)}, \VUgras{mitxoletak} \textsuperscript{(2)} eta
\VUgras{kukupraka} \textsuperscript{(3)} eztabaida alaia sartzen dute galtzen
tonu berdinean beren \VUgras{koroa} freskoekin.

%%K t08-c1-03-1 p
3. --- Fruta-arbolak hostoz jantzi dira; fruituak heldu dira. \VUgras{Gerezondo}
txikia \textsuperscript{(4)} makurtzen da \VUgras{gerezi} ederren pean
\textsuperscript{(5)}, zeinak haurrek gogotik biltzen eta jaten baitituzte.

%%K t08-c1-04-1 p
4. --- Artaldea egun osoan egon daiteke orain aske.

%%K t08-c1-04-2 p
\VUgras{Bildotsak} \textsuperscript{(6)} hazi dira eta \VUgras{ardi} eder
\textsuperscript{(7)} eta \VUgras{ahari} eder \textsuperscript{(8)} bihurtu dira,
artile lodikoak. Lasai jaten dute \VUgras{belarra} \VUgras{larrean}
\textsuperscript{(9)}, \VUgras{bitxilorez} josita.

%%K t08-c1-04-3 p
Laster moztuko diete animalia hauei \VUgras{artilea}, zeinaz gure arropetarako
oihala egiten baita.

%%K t08-tit-2 sub
\VUsaut{5.0mm}
\VUpk{10.2pt}{40}{Artzaina.}
\VUsaut{3.4mm}

%%K t08-c1-05-1 p
5. --- \VUgras{Artzainak} \textsuperscript{(10)}, dirudienez, ez die askorik
begiratzen. Zeruertzean igarotzen ari den ekaitzak bakarrik kezkatzen du, eta
\VUgras{artzain-mutila} \textsuperscript{(13)}, bere \VUgras{zahagia}
\textsuperscript{(14)} ahora eramanez, are axolagabeagoa dirudi.

%%K t08-c1-06-1 p
6. --- Beroak zapaltzen du; hara non \VUgras{txakurra} \textsuperscript{(15)}
freskatzera doan; \VUgras{errekako} \textsuperscript{(16)} ur freskoa miazkatzen
du eta \VUgras{igel} txiki gaixoa \textsuperscript{(17)} izutzen, zeina ertzeko
belarrean ezkutatuta baitzegoen. Hark ur gardenera egiten du jauzi, non
\VUgras{igebelarra} \textsuperscript{(18)} loratzen baita.

%%K t08-c1-07-1 p
7. --- \VUgras{Buztanikarak} \textsuperscript{(19)} bainua hartzen du
\VUgras{iturburu} txikiaren ondoan \textsuperscript{(20)}, zeina bi harriren
artean sortzen baita eta \VUgras{sasi arantzatsuen} \textsuperscript{(21)} eta
\VUgras{elorri-sastraken} \textsuperscript{(22)} azpian ezkutatzen.

%%K t08-tit-3 sub
\VUsaut{5.0mm}
\VUpk{10.2pt}{40}{Uzta.}
\VUsaut{3.4mm}

%%K t08-c1-08-1 p
8. --- Herriko haurrentzat garia biltzea garai oso alaia da. Pozik
\VUgras{itzulipurdika} \textsuperscript{(24)} aritzen dira \VUgras{lasto-metetan}
\textsuperscript{(23)}, \VUgras{igitariei} \textsuperscript{(26)} laguntzen
diete, eta haiek \VUgras{soroa} \textsuperscript{(27)} beren \VUgras{sega}
\textsuperscript{(25)} ondo zorroztuekin ebakitzen duten bitartean ---
\VUgras{harri-zorroztailean} \textsuperscript{(30)} zorroztuak --- haurrek garia
biltzen dute eta \VUgras{azao} \textsuperscript{(31)} edo \VUgras{sorta}
\textsuperscript{(32)} egiten.

%%K t08-c1-09-1 p
9. --- Batzuetan handienei uzten diete \VUgras{igitaiaz} \textsuperscript{(12)}
\VUgras{urrezko zurtoinak} \textsuperscript{(28)} ebakitzen, zeinek
\VUgras{galburu} astunak \textsuperscript{(29)} daramatzaten, alez beteak; eta
txikiek, \VUgras{abereak} \textsuperscript{(33)} zainduz --- \VUgras{belardian}
\textsuperscript{(34)} bazkatzen baitira \VUgras{ezki} handiaren
\textsuperscript{(35)} itzalpean ---, \VUgras{tximeleta} ederrak harrapatzen
dituzte beren \VUgras{sarearekin} \textsuperscript{(36)}.

%%K t08-c1-09-2 p
Batzuk \VUgras{gurdira} \textsuperscript{(37)} ere igotzen dira, \VUgras{sardez}
\textsuperscript{(38)} ematen dizkieten azaoak jartzeko.

%%K t08-c1-10-1 p
10. --- \VUgras{Lautada} osoan \textsuperscript{(39)}, non \VUgras{hegaberak}
\textsuperscript{(41)} hegaldatzen baitira, iskanbila handia dago. Denak uztan ari
dira. Baina behartsuek lehen bezala tresna zaharrei heltzen dieten bitartean ---
\VUgras{segak, igitaiak} eta \VUgras{iharrak} \textsuperscript{(40)} --- lurjabe
aberatsek beren garia edo beren \VUgras{zekalea} \textsuperscript{(43)}
\VUgras{uzta-makinaz} \textsuperscript{(42)} biltzen dute. Gero, azaoak lastategira
eraman gabe, bertan bertan \VUgras{meta} handiak \textsuperscript{(44)} egiten
dituzte.

%%K t08-c1-11-1 p
11. --- Gero \VUgras{jotzeko makina} \textsuperscript{(47)} ekartzen dute, zeina
\VUgras{lokomobilak} \textsuperscript{(45)} biratzen baitu \VUgras{transmisio-uhalaren}
\textsuperscript{(46)} bidez, eta hala alea \VUgras{lastotik} bereizten da, erein
zuten landa utzi gabe.

%%K t08-tit-4 sub
\VUsaut{5.0mm}
\VUpk{10.2pt}{40}{Ekaitza.}
\VUsaut{3.4mm}

%%K t08-c1-12-1 p
12. --- Uzta garaian sarritan \VUgras{ekaitz} handiak \textsuperscript{(53)}
lehertzen dira. Lehenik urrutitik \VUgras{trumoiaren} danbatekoa entzuten da;
\VUgras{mendebaleko haize} indartsua \textsuperscript{(54)} altxatzen da eta
\VUgras{basoko} \textsuperscript{(49)} zuhaitzik handienak intzirika makurrarazten
ditu. \VUgras{Hodei} beltzek \textsuperscript{(56)} zerua hartzen dute
erasoka; \VUgras{tximistaren} \textsuperscript{(55)} sigi-saga itsugarriek
zeharkatzen dituzte. \VUgras{Euria} hasten da erortzen, eta batzuetan
\VUgras{txingorra} ere, uztarako prest dagoen garia jota.

%%K t08-c1-13-1 p
13. --- Sarritan tximistak eraikinen bat pizten du. Hala, iaz \VUgras{haize-errotaren}
\textsuperscript{(50)} teilatua erabat erre zen, eta haren bi \VUgras{hegal}
\textsuperscript{(51)} hautsi zituzten.

%%K t08-c1-14-1 p
14. --- Ordu batzuen buruan bakea itzultzen da. Zeruertzean \VUgras{ostadar}
distiratsua \textsuperscript{(52)} agertzen da, eta naturak berriro hartzen du
denboraldi batez etenda zegoen bizitza. \VUgras{Txaboletan}
\textsuperscript{(48)} babestutako herritarrak lanera itzultzen dira eta ausarki
hasten dira jasan berri duten kaltea konpontzen.

%%K t08-tit-5 sub
\VUsaut{6.0mm}
\VUcentre{13.2pt}{90}{\textit{Bigarren agerraldia.}}
\VUsaut{3.6mm}

%%K t08-tit-6 sub
\VUpk{10.2pt}{40}{Ehiza. --- Mahats-bilketa.}
\VUsaut{1.4mm}
\VUpk{10.2pt}{40}{Arrantza.}
\VUsaut{4.0mm}

%%K t08-c2-01-1 p
1. --- \VUgras{Abuztuaren} amaieran edo \VUgras{irailaren} erdialdean, eskualdearen
arabera, ehiza-denboraldia irekitzen da. Eguzkiaren lehen izpiek
\VUgras{musker} \textsuperscript{(2)} zuloetatik ateratzen dituzten orduko,
\VUgras{ehiztaria} \textsuperscript{(5)} irteten da jada bere \VUgras{txakurrekin}
\textsuperscript{(4)}.

%%K t08-c2-02-1 p
2. --- Bere \VUgras{ehiza-jantzi} sendoa \textsuperscript{(9)} jantzi du, oihal
lodizkoa, eta larruzko \VUgras{polainak}, zeinek belar bustian ibiltzen utziko
baitiote, non \VUgras{apoak} \textsuperscript{(1)} ezkutatzen baitira;
\VUgras{eskopeta} \textsuperscript{(6)} eta \VUgras{ehiza-zorroa}
\textsuperscript{(10)} hartu ditu --- eta bidera.

%%K t08-c2-03-1 p
3. --- Ehizaren sugarrak mahastiaren erdi-erdira eramaten du, non bilketa bete-betean
baitago. Mahastizainaren haurrek, zeinek normalean ahuntzak zaintzen baitituzte
bidearen ondoan, gaur nahi duten lekuan uzten dituzte bazkatzen, eta \VUgras{aker}
zaharra \textsuperscript{(3)} hesola bati lotuta --- hark seguru asko askatasuna
baliatuko bailuke ihes egiteko ---, beren poz-partea hartzen dute ere. Batek
mordo bat ateratzen du \VUgras{mahats-saskitik} \textsuperscript{(12)} eta
\VUgras{zukua} \textsuperscript{(14)} \VUgras{katilura} \textsuperscript{(13)}
estutzen olgatzen da, muztioa (hartzitu gabeko ardoa) atera dadin. Besteak
\VUgras{hosto-koroaz} \textsuperscript{(15)} koroatzen du bere burua, josi berri
duenaz. Haien ondoan \VUgras{txolarre} ausartak \textsuperscript{(16)} dolarearen
ondo-ondora hurbiltzen dira, mahatsa ostera.

%%K t08-c2-04-1 p
4. --- Haurrak olgatzen diren bitartean, gizonak lanean ari dira.
\VUgras{Biltzaileek} \textsuperscript{(18)} mahatsa sotora eramaten dute
\VUgras{bizkar-saskietan} \textsuperscript{(17)}; \VUgras{kupelgileak}
\textsuperscript{(19)} \VUgras{kupelak} \textsuperscript{(20)} konpontzen ditu,
\VUgras{oholak} \textsuperscript{(22)} eta \VUgras{hondoak}
\textsuperscript{(21)} osorik dauden begiratzen du, eta hautsitako
\VUgras{uztaiak} \textsuperscript{(23)} ordezkatzen. Berak egin zituen
\VUgras{upel} handiak ere \textsuperscript{(25)}, non \VUgras{mahatsa}
\textsuperscript{(26)} isurtzen baita hartzitu dadin.

%%K t08-c2-05-1 p
5. --- Hartzidura amaitzen denean, ardoa isurtzen da, eta hondarra
\VUgras{dolarearen} \textsuperscript{(28)} azpian jartzen. \VUgras{Torloju}
handia \textsuperscript{(29)} pixkanaka estutuz, zukua ateratzen da, zeina
\VUgras{upelera} \textsuperscript{(32)} isurtzen baita eta oraindik \VUgras{ardo}
gehiago \textsuperscript{(31)} ematen.

%%K t08-tit-8 sub
\VUsaut{6.0mm}
\VUpk{10.2pt}{40}{Garagardoa eta sagardoa.}
\VUsaut{3.4mm}

%%K t08-c2-06-1 p
6. --- \VUgras{Upategiaren} \textsuperscript{(27)} hormatik gora
\VUgras{lupuluaren} \textsuperscript{(30)} adar bat kiribiltzen da. Hemen apaingarri
hutsa da, baina beste herrialde batzuetan, non mahatsondoa oso urria baita,
lupulua lantzen dute \VUgras{garagardoa} egiteko.

%%K t08-c2-07-1 p
7. --- \VUgras{Sagarrondo} txikia \textsuperscript{(33)} sagarren pean makurtzen
da, zeinek \VUgras{sagardo} bikaina emango baitute helduko direnean. Beren
\VUgras{kaiolan} \textsuperscript{(34)} \VUgras{kanarioak}
\textsuperscript{(35)} eta \VUgras{karnabak} \textsuperscript{(36)} bekaizkeriaz
begiratzen diete aske dabiltzan txolarreei.

%%K t08-c2-08-1 p
8. --- Begiak iristen diren guztian, muinoen hegalean \VUgras{hesola}
\textsuperscript{(40)} lerroak daude, \VUgras{mahatsondo} bikainak
\textsuperscript{(39)} eusten dituztenak, \VUgras{mordoez} astunduak.

%%K t08-c2-08-2 p
Urte osoan \VUgras{mahastizainak} \textsuperscript{(42)} lan egin du bere
\VUgras{mahastia} \textsuperscript{(38)} zaintzeko, eta hara non saritua den
uzta onarekin. \VUgras{Biltzaileek} \textsuperscript{(41)} gurdi gainean
jarritako upelak betetzen dituzte, \VUgras{mandoek} \textsuperscript{(43)}
tiratzen dutena, eta denak alai daude.

%%K t08-tit-9 sub
\VUsaut{5.0mm}
\VUpk{10.2pt}{40}{Arrantza.}
\VUsaut{3.4mm}

%%K t08-c2-09-1 p
9. --- Gauza bera \VUgras{arrantzaleentzat} \textsuperscript{(44)}, zeinak
egunsentian etorri baitira uraren ondoan kokatzera beren \VUgras{kanaberekin}
\textsuperscript{(45)}.

%%K t08-c2-09-2 p
\VUgras{Arrainak} \textsuperscript{(47)} amua hartzen du beren \VUgras{haria}
\textsuperscript{(46)} uretan sartu orduko, eta ozta-ozta dute astirik
\VUgras{beita} aldatzeko, biktima berri bat hariaren muturrean saltoka baitago
jada.

%%K t08-c2-09-3 p
Hala ere, \VUgras{sarezko arrantzaleak} \textsuperscript{(48)}, zeinak bere sarea
\VUgras{txaluperik} \textsuperscript{(49)} \VUgras{ibaiaren}
\textsuperscript{(50)} erdira botatzen baitu, askoz gehiago harrapatzen du.

%%K t08-c2-10-1 p
10. --- Udazkena ibilaldi onen garaia da: oinez, \VUgras{zaldiz}
\textsuperscript{(52)}, \VUgras{bizikletaz} \textsuperscript{(53)},
\VUgras{automobilez} \textsuperscript{(54)}. \VUgras{Bideak}
\textsuperscript{(51)} garbi daude, eta jendeak azken egun eguzkitsuak baliatzen
ditu, \VUgras{enarak} \textsuperscript{(56)} teilatuetan biltzen ari baitira jada
herrialde epeletara hegan egiteko. \VUgras{Hurritz} zaharraren
\textsuperscript{(57)} hostoak horitzen ari dira, \VUgras{hurrak} erortzen, eta
\VUgras{basotxo} \textsuperscript{(58)} bat osatzen duten \VUgras{zuhaitz}
garaiak, \VUgras{goragunekoak} \textsuperscript{(59)}, laster biluziko dira.

%%K t08-c2-11-1 p
11. --- \VUgras{Eguzkia} \textsuperscript{(60)} beranduago jaikitzen da, egunak
laburtzen ari dira, \VUgras{goizean} hotz nahiko zorrotza sumatzen hasten da, eta
\VUgras{oilagor} \textsuperscript{(61)} taldeak gure herrialde etsaia uzten ari
dira jada.
\VUsaut{6.0mm}
\VUfilet{22mm}
